% Generated draft: requires mathematical and source review.
\documentclass[11pt]{article}
\usepackage[margin=25mm]{geometry}
\usepackage{fontspec}
\setmainfont{Latin Modern Roman}
\usepackage{amsmath,amssymb,mathtools,unicode-math}
\setmathfont{Latin Modern Math}
\usepackage{enumitem,needspace,xurl,hyperref}
\setlength{\parindent}{0pt}
\setlength{\parskip}{5pt}
\setlength{\emergencystretch}{3em}
\newcommand{\N}{\mathbb{N}}
\newcommand{\Z}{\mathbb{Z}}
\newcommand{\Q}{\mathbb{Q}}
\newcommand{\R}{\mathbb{R}}
\newcommand{\C}{\mathbb{C}}
\newcommand{\F}{\mathbb{F}}
\newcommand{\A}{\mathbb{A}}
\newcommand{\PP}{\mathbb{P}}
\newcommand{\E}{\mathbb{E}}
\newcommand{\eps}{\varepsilon}
\newcommand{\dd}{\,\mathrm{d}}
\begin{document}
\section*{UnsolvedMath ID: 30000330}
\textbf{Source problem code:} catalogue entry
\subsection{Definitions and mathematical statement}
In the usual complete connected interpretation, a hyperbolic surface $S$ is a complete connected Riemannian two-manifold of constant curvature $-1$. It is $K$-quasiconformally homogeneous if
\[
 \forall x,y\in S\ \exists f:S\to S\text{ a quasiconformal homeomorphism},
 \qquad f(x)=y,\quad K(f)\le K,
\]
where $K(f)$ is maximal quasiconformal dilatation. One common constant $K$ must work for every pair; the maps may vary with the pair. No fixed transitive group all of whose elements have dilatation at most $K$ is required.

The dimension-two uniform-gap question is
\[
 \exists K_2>1\ \forall S\not\cong\mathbb H^2\ \forall K\ge1,
 \qquad S\text{ is }K\text{-quasiconformally homogeneous}
 \Longrightarrow K\ge K_2.
\]
The exclusion is of the hyperbolic plane as a hyperbolic surface. The constant is universal, independent of the surface's genus, area or end structure. The source does not restrict to compact, finite-area or finite-topological-type surfaces. Showing only that each individual nonplane surface requires some constant greater than one would be weaker than the uniform assertion.

Completeness is part of the standard formulation used here; the extract does not separately state it or specify orientability and whether the maps must preserve orientation. Those conventions should match the report rather than adding a finite-type restriction.

The literature assessment checked on 21 August 2026 cites only the original report and repeats its statement that the two-dimensional analogue was unknown. This preserves the historical source status but is not an independently established current-status conclusion.
\subsection{Short English statement}
Is there one universal gap above dilatation one for every quasiconformally homogeneous hyperbolic surface except the hyperbolic plane? The bound must be independent of the surface, with no compactness or finite-type restriction added.
\subsection{Sources}
\begin{itemize}
\item \emph{Low-Dimensional Manifolds} (2006), Oberwolfach Reports, DOI 10.4171/owr/2005/45, pp. 2519--2570, \url{https://doi.org/10.4171/owr/2005/45}; proposed year supplied as 2005.
\item The sole source in the assessment checked 21 August 2026 is the original report, \url{https://doi.org/10.4171/owr/2005/45}; its dimension-two unknown-status statement is historical evidence.
\item Local UnsolvedMath registry, permanent ID $30000330$, exact source code \texttt{catalogue entry}, created 25 August 2026 and updated 16 September 2026. Supplied verification restores the full uniform bound in place of a comparison reference.
\end{itemize}
\end{document}
